% =====================================================================
%  5.4 Cosinussetningen — fasit
% =====================================================================

\fk{Del 1 — Cosinussetningen og beviset}
\fo{1.1} a) $10$ \ b) $2\sqrt7$ \ c) $\sqrt{19}$
\fo{1.2} $\cos A = \frac{6^2 + 16^2 - 14^2}{2\cdot6\cdot16} = \frac{96}{192} = \frac12$, så $A = 60^\circ$
\fo{1.3} a) $a \approx 11{,}3$ cm \ b) $\cos^{-1}\!\left(\frac{6^2 + 8^2 - 11^2}{2\cdot 6\cdot 8}\right) = \cos^{-1}\!\left(-\frac{21}{96}\right) \approx 102{,}6^\circ$ (vinkelen overfor $11$ cm)
\fo{1.4} a) $\cos A = \frac{2^2 + 2^2 - (2\sqrt3)^2}{2\cdot2\cdot2} = -\frac12$, så $120^\circ$ \ b) $\cos A = \frac{2+2-4}{4} = 0$, så $90^\circ$
\fo{1.5} a) $PQ^2 = 420^2 + 310^2 - 2\cdot420\cdot310\cdot\cos 72^\circ$, så $PQ \approx 438{,}215$ m $\approx 438{,}2$ m \ b) Vinkelen ved $P$ står overfor $RQ = 310$ m: $\cos P = \frac{420^2 + 438{,}215^2 - 310^2}{2\cdot 420\cdot 438{,}215}$, så $P \approx 42{,}283^\circ \approx 42{,}3^\circ$
\fo{1.6} Feilen er overgangen til $a^2 = 136 - 60$. Det er $-120\cdot(-\frac12) = +60$ (minus ganger minus), så $a^2 = 136 + 60 = 196$ og $a = 14$.
\fo{1.7} a) $x = 8\cos 60^\circ = 4$ og $h = \sqrt{8^2 - 4^2} = 4\sqrt3$ \ b) $a^2 = h^2 + (10-4)^2 = 48 + 36 = 84$ \ c) $a^2 = 64 + 100 - 160\cdot\frac12 = 84$. Samme svar.
\fo{1.8} Ligger $B$ som stump vinkel, faller høyden fra $C$ ned på forlengelsen av $c$, forbi $B$. Da er $D$ ikke mellom $A$ og $B$, og $DB = x - c$.
Pytagoras i trekant $DBC$ gir $a^2 = h^2 + (x-c)^2$. Siden $(x-c)^2 = (c-x)^2$, er regningen den samme som i beviset: $a^2 = h^2 + c^2 - 2cx + x^2 = b^2 + c^2 - 2cx$, og med $x = b\cos A$ får vi $a^2 = b^2 + c^2 - 2bc\cos A$.
\fo{1.9} a) Alltid. Fordi $\cos A > -1$ for $0^\circ < A < 180^\circ$ er $a^2 < b^2 + c^2 + 2bc = (b+c)^2$, så $a < b + c$.
b) Noen ganger. $A = 90^\circ$, $b = 3$, $c = 4$ gir $a = 5$, som er størst. Med samme $b = 3$ og $c = 4$ gir $A = 30^\circ$ at $a^2 = 25 - 12\sqrt3 \approx 4{,}2$, så $a \approx 2{,}0$, som er kortest.
c) Alltid. Da er $\cos A = \frac{b^2 + c^2 - a^2}{2bc} = \frac{bc}{2bc} = \frac12$. Mellom $0^\circ$ og $180^\circ$ finnes bare én vinkel med cosinus $\frac12$, nemlig $60^\circ$.
\fo{1.10} $(x+2)^2 = x^2 + (x+1)^2 - 2x(x+1)\cdot\frac18$. Ganger med $4$ og samler: $3x^2 - 9x - 12 = 0$, altså $x^2 - 3x - 4 = 0$. Løsningene er $x = 4$ og $x = -1$. Siden $x > 0$, er $x = 4$ (sidene $4$, $5$, $6$).
\fo{Sjekk 1} a) $a = 7$ \ b) $\cos A = \frac{b^2 + c^2 - a^2}{2bc}$ \ c) Når $A$ er stump, er $\cos A < 0$, så $-2bc\cos A > 0$ og $a^2 = b^2 + c^2 + (\text{noe positivt})$.

\fk{Del 2 — Velg riktig setning}
\fo{2.1} a) SSS, ingen kjent vinkel–side-par: cosinussetningen \ b) SVS, vinkelen ligger mellom sidene og det finnes ikke noe vinkel–side-par: cosinussetningen \ c) SVV, $A$ og $a$ er et par: sinussetningen \ d) SSV, $B$ og $b$ er et par: sinussetningen. Cosinussetningen går også, men gir en andregradslikning.
\fo{2.2} a) SVS, ingen kjent vinkel–side-par: cosinussetningen, $a \approx 13{,}1$ \ b) SVV, $A$ og $a$ er et par: sinussetningen, $b \approx 12{,}2$ \ c) SSS: cosinussetningen, $C = \cos^{-1}\!\left(\frac{8^2 + 11^2 - 13^2}{2\cdot 8\cdot 11}\right) = \cos^{-1}\!\left(\frac{1}{11}\right) \approx 84{,}8^\circ$
\fo{2.3} a) Vinkelsummen stemmer, men den lengste siden $b = 7$ ligger overfor $60^\circ$, ikke overfor den største vinkelen $80^\circ$. \ b) Vinkelsummen er $190^\circ$, ikke $180^\circ$.
\fo{2.4} a) $169 = 225 + c^2 - 15c$ gir $c^2 - 15c + 56 = 0$, så $c = 7$ eller $c = 8$. To trekanter.
b) $361 = 25 + c^2 - 5c$ gir $c^2 - 5c - 336 = 0$, så $c = 21$ eller $c = -16$. Én trekant.
c) $36 = 64 + c^2 - 8c$ gir $c^2 - 8c + 28 = 0$ med $64 - 112 < 0$. Ingen trekant.
\fo{2.5} Vinklene står overfor sidene $7$, $9$ og $12$. $C = \cos^{-1}\!\left(\frac{7^2 + 9^2 - 12^2}{2\cdot 7\cdot 9}\right) = \cos^{-1}(-\frac19) \approx 96{,}379^\circ \approx 96{,}4^\circ$. $B = \cos^{-1}\!\left(\frac{7^2 + 12^2 - 9^2}{2\cdot 7\cdot 12}\right) = \cos^{-1}(\frac23) \approx 48{,}190^\circ \approx 48{,}2^\circ$.
Sinussetningen: $\sin A = \frac{7\sin C}{12} \approx 0{,}580$ gir $A \approx 35{,}431^\circ \approx 35{,}4^\circ$. Kontroll: $96{,}4^\circ + 48{,}2^\circ + 35{,}4^\circ = 180^\circ$.
Sinus er trygt for $A$ fordi $C$ er den eneste vinkelen som kan være stump, og $A$ og $B$ er derfor spisse. Da gir $\sin A$ bare én mulig vinkel.
\fo{2.6} Første feil er $C \approx 56{,}8^\circ$. Likningen $\sin C \approx 0{,}837$ har også løsningen $C \approx 123{,}2^\circ$. Siden $c = 9$ er den lengste siden, må $C$ være den største vinkelen. Med $56{,}8^\circ$ blir $B = 107{,}0^\circ$ størst, selv om $b < c$, og det går ikke.
Riktig: $C = \cos^{-1}\!\left(\frac{3^2 + 7^2 - 9^2}{2\cdot 3\cdot 7}\right) = \cos^{-1}(-\frac{23}{42}) \approx 123{,}204^\circ \approx 123{,}2^\circ$ og $B = \cos^{-1}(\frac{41}{54}) \approx 40{,}601^\circ \approx 40{,}6^\circ$. Sum: $16{,}2^\circ + 123{,}2^\circ + 40{,}6^\circ = 180^\circ$.
\fo{2.7} $\cos C = \frac{9^2 + 10^2 - 11^2}{2\cdot9\cdot10} = \frac13$, så $C \approx 70{,}529^\circ$ (mellom $9$ og $10$). $T = \frac12\cdot9\cdot10\cdot\sin 70{,}529^\circ \approx 42{,}4$ m$^2$.
\fo{2.8} a) Alltid. Da er $b^2 + c^2 - a^2 < 0$ og $2bc > 0$, så $\cos A < 0$ og $A$ er stump.
b) Noen ganger. $2, 3, 4$: $4 + 9 < 16$, stump. $4, 4, 5$: $16 + 16 > 25$, så $\cos A > 0$ og $A$ er spiss.
c) Aldri. Cosinus til en vinkel er alltid minst $-1$, så en verdi under $-1$ betyr at vinkelen ikke finnes. Eksempel: sidene $4$, $6$ og $11$ gir $\cos A = \frac{4^2 + 6^2 - 11^2}{2\cdot4\cdot6} = -\frac{23}{16} < -1$ (også $4 + 6 < 11$).
\fo{2.9} Cosinussetningen: $a^2 = 100 + c^2 - 10c$, altså $c^2 - 10c + 100 - a^2 = 0$ med røttene $c = 5 \pm \sqrt{a^2 - 75}$.
Ingen røtter når $a^2 < 75$, altså $a < 5\sqrt3$: \textbf{ingen trekant}.
Én løsning når $a = 5\sqrt3$ ($c = 5$). Én \emph{positiv} rot når $a \ge 10$: for $a > 10$ er den andre roten negativ, og $a = 10$ gir røttene $0$ og $10$. \textbf{Én trekant}.
To positive røtter når $5 - \sqrt{a^2 - 75} > 0$, altså $a^2 < 100$. Så $5\sqrt3 < a < 10$: \textbf{to trekanter}. Grensene er de samme som i 5.3 med sinussetningen: $b\sin A = 5\sqrt3$ og $b = 10$.
\fo{2.10} Likningen $a^2 = b^2 + c^2 - 2bc\cos A$ gir $c^2 - 2b\cos A\cdot c + b^2 - a^2 = 0$, så $c = b\cos A \pm \sqrt{a^2 - b^2\sin^2 A}$.
Når $A$ er stump, er $\cos A < 0$. Da er $b\cos A - \sqrt{\ldots} < 0$, så minus-roten er negativ og gir ingen trekant. Det er derfor høyst én trekant.
Pluss-roten er positiv nøyaktig når $\sqrt{a^2 - b^2\sin^2 A} > -b\cos A$, altså $a^2 > b^2\sin^2 A + b^2\cos^2 A = b^2$, så $a > b$. Da finnes én trekant. Er $a \le b$, finnes ingen (den stumpe vinkelen må ligge overfor den lengste siden).
\fo{Sjekk 2} a) Sinussetningen (SVV) \ b) Cosinussetningen (SSS) \ c) $\cos v$ gir bare én vinkel mellom $0^\circ$ og $180^\circ$.

\fk{Ekstra utfordring}
\fo{U1} a) \lstinline[style=laereverk,numbers=none,frame=none]|abs(storst - 90) < 0.0001|. Flyttall er avrundet. For sidene $1, 1, \sqrt2$ blir vinkelen $90{,}00000000000001$ i stedet for $90$, så \lstinline[style=laereverk,numbers=none,frame=none]|storst == 90| ville ha gitt «stump».
b) $3,4,5$: \texttt{36.9 53.1 90.0} og «rett». $5,7,8$: \texttt{38.2 60.0 81.8} og «spiss». $1,2,4$: «Ingen trekant». $1,1,\sqrt2$: \texttt{45.0 45.0 90.0} og «rett».
c) For $1,2,4$ ville cosinus bli utenfor $-1$ til $1$, og \lstinline[style=laereverk,numbers=none,frame=none]|acos| gir feilmelding. Ulikheten fanger dette først.
\fo{U2} a) $6,7,10$: $36 + 49 = 85 < 100$, stump \ b) $8,9,12$: $64 + 81 = 145 > 144$, spiss \ c) $20,21,29$: $400 + 441 = 841 = 29^2$, rett.
Forklaring: Fortegnet til $b^2 + c^2 - a^2$ er fortegnet til $\cos A$. En trekant har høyst én vinkel på $90^\circ$ eller mer, og den ligger overfor den lengste siden. Er den vinkelen spiss, er de andre det også.
\fo{U3} a) $B = \cos^{-1}\!\left(\frac{35^2 + 48^2 - 60^2}{2\cdot35\cdot48}\right) \approx 91{,}211^\circ \approx 91{,}2^\circ$ og $D = \cos^{-1}\!\left(\frac{42^2 + 52^2 - 60^2}{2\cdot42\cdot52}\right) \approx 78{,}538^\circ \approx 78{,}5^\circ$
b) $T \approx \frac12\cdot35\cdot48\cdot\sin 91{,}2^\circ + \frac12\cdot42\cdot52\cdot\sin 78{,}5^\circ \approx 840 + 1070 \approx 1910$ m$^2$
\fo{U4} Trekant $ABD$ har vinkel $v$ ved $A$: $BD^2 = a^2 + b^2 - 2ab\cos v$. Vinkelen ved $B$ er $180^\circ - v$, så i trekant $ABC$ er $AC^2 = a^2 + b^2 - 2ab\cos(180^\circ - v) = a^2 + b^2 + 2ab\cos v$. Summen er $AC^2 + BD^2 = 2a^2 + 2b^2 = a^2 + b^2 + a^2 + b^2$.
\fo{U5} Ved $120^\circ$ er $a^2 = b^2 + c^2 + bc$. Prøv par $b, c$ til høyre side er et kvadrattall. $b = 7$, $c = 8$ gir $49 + 56 + 64 = 169$, så $a = 13$. Kontroll: $\cos A = \frac{7^2 + 8^2 - 13^2}{2\cdot7\cdot8} = \frac{-56}{112} = -\frac12$. Flere svar er riktige, for eksempel $5, 16, 19$.
